\section{Orbit-finitely spanned vector spaces}\label{sec:spaces}
We now introduce the main topic of this paper: vector spaces with an orbit-finite spanning set. 
We begin with the special case of spaces that have an orbit-finite basis; 
this special case has an elementary definition, and yet it will be the relevant case for almost all results of this paper. 

\begin{definition}[Orbit-finite-dimensional vector space]\label{def:orbit-finite-basis}
    $\Lin_\field X$ is --- given an orbit-finite set $X$ and a field $\field$ --- the vector space of finite formal linear combinations of elements in $X$ over $\field$.
\end{definition}

This definition has two important parameters: the field $\field$ and the oligomorphic structure $\A$ over which $X$ is an orbit-finite set.
The spaces defined here have two kinds of structure: that of a vector space, and the action of the automorphism group of $\A$
--- via $\pi(\sum_i \lambda_i x_i) = \sum_i \lambda_i \pi(x_i)$. 
We will be interested in subsets that preserve both kinds of structure, 
i.e.~those that are closed under taking linear combinations as well as applying automorphisms of $\A$. 
Such subsets are called \emph{equivariant subspaces}. 


\begin{myexample}\label{ex:length-one-dim-one}
    Let $\A$ be the equality atoms and let $\field$ be an arbitrary field. 
    As explained in~\cite[Ex.~4.2]{BFKM24}, the corresponding vector space $\Lin_\field \A$ has only three equivariant subspaces: 
    the zero subspace, the full space, and the subspace consisting of vectors whose coefficients add up to zero 
    (e.g.~$a - b$ and $a + b - 2 c$, with $a, b, c \in \A$ distinct). 
\end{myexample}

 Unfortunately, there is a price to pay for the elementary character of Definition~\ref{def:orbit-finite-basis}, 
 and that is the failure of certain closure properties. 
 In particular, these spaces are not closed --- even up to isomorphisms, i.e.~equivariant linear bijections --- under taking equivariant subspaces or quotients under such spaces. 
 The problem with subspaces is apparent already in Example~\ref{ex:length-one-dim-one}, since the unique nontrivial subspace of $\Lin_\field \A$ does not have any equivariant basis, regardless of the underlying field~\cite[Ex.~6.1]{BFKM24}; 
 this subspace with no orbit-finite basis can also be realised as a quotient of $\Lin_\field \A^2$.
 % (Also an orbit-finite-dimensional vector space can have two orbit-finite bases without an equivariant bijection between them \cite[\S3.4]{Ghosh25}.)
 For this reason, we use a more general notion of vector spaces, which is presented below, in a style that emphasises the group action, as was done in Definition~\ref{def:orbit-finite-set-alternative}.

\begin{definition}[Orbit-finitely spanned vector space, abstractly]\label{def:orbit-finitely-spanned-abstract}
    An \emph{orbit-finitely spanned vector space} over an oligomorphic structure $\A$ is a vector space equipped with an action of automorphisms of $\A$ such that: 
    \begin{enumerate}
        \item \label{item:equivariance-of-vector-space-structure} 
        vector addition and scalar multiplication are equivariant, with the action on the field being trivial;%
        \footnote{Equivalently, the group action $v \mapsto \pi(v)$ of each $\pi$ is linear.}

        \item 
        every vector is supported by a finite set of atoms;%
        \footnote{Equivalently, the application $(\pi, v) \mapsto \pi(v)$ is continuous, where the automorphism group is endowed with the topology of pointwise convergence whilst the vector space is endowed with the discrete topology. 
        See \cite[Lemma~4.1.5]{hodges1993model}.}

        \item \label{item:orbit-finite-spanning-subset} 
        the space is spanned by some subset that is orbit-finite.%
        \footnote{Equivalently, the vector space is finitely generated as a module over the group ring of automorphisms.}
    \end{enumerate}
\end{definition}

The spaces defined above are easily seen to be closed under taking quotients by equivariant subspaces. 
Closure under taking equivariant subspaces is less obvious: 
one could imagine that condition~\eqref{item:orbit-finite-spanning-subset} above is violated by moving to an equivariant subspace. 
This closure property turns out to be intimately related to the ascending chain property discussed in the introduction:

\begin{theorem}\label{thm:ACP}
    For any field $\field$ and oligomorphic structure $\A$, the following conditions are equivalent: 
    \begin{enumerate}
        \item orbit-finitely spanned vector spaces are closed under taking equivariant subspaces; 
        \item for every $d \in \set{1,2,\ldots}$,  the vector space 
        $
            \Lin_\field \A^d
        $ does not have any infinite ascending chain of equivariant subspaces;
        \item for every orbit-finitely spanned vector space $V$, there are no infinite ascending chains of equivariant subspaces in $V$.
    \end{enumerate}
    Furthermore, if these conditions hold, then every orbit-finitely spanned vector space is isomorphic (via an equivariant linear bijection) to one of the form $U/W$, 
    where $W \subseteq U$ are equivariant subspaces of $\Lin_\field \A^d$ for some $d \in \set{0,1,\ldots}$.
\end{theorem}


We say that an atom structure $\A$ has the \emph{ascending chain property} over a field $\field$ if any of the equivalent conditions in Theorem~\ref{thm:ACP} is satisfied. 
One interpretation of the theorem is that the ascending chain property is necessary for the theory of vector spaces to be well-behaved. 
In particular, thanks to the ``furthermore'' part, we get a similar result to the equivalence of Definitions~\ref{def:orbit-finite-set} and~\ref{def:orbit-finite-set-alternative} in the previous section, 
i.e.~a concrete characterisation of the orbit-finitely spanned vector spaces that can be used in algorithms (provided we know how to represent equivariant subspaces --- see Section~\ref{sec:equivariant-subspaces}). 


We are therefore interested in atom structures that have the ascending chain property. 
As it turns out, the techniques used in this paper will yield a stronger property, namely a finite bound on the length of chains: 




\begin{definition}
    [Finite length property]
    \label{def:finite-length-property}
    An oligomorphic structure $\A$ has the \emph{finite length property over a field $\field$} if for every orbit-finite set $X$ over $\A$, there is a finite upper bound on the length of chains of equivariant subspaces 
    of $\Lin_\field X$.
    (The supremum of the chain lengths, finite or not, is called the \emph{length} of $\Lin_\FF X$.)
\end{definition}

In light of Theorem~\ref{thm:ACP} and the proof of (2) $\Rightarrow$ (3) in Appendix~\ref{sec:appendix-modules}, we could have used $\A^d$ instead of $X$ in the above definition.
Like mentioned in the introduction, the finite length property can be strictly stronger than the ascending chain property, as witnessed by the vector atoms from Example~\ref{ex:vector-space-atoms}. 
(Note that when we talk about the vector atoms, there are two fields involved, namely the finite field used to define $\A$ and the field $\field$ used to define $\Lin_\field \A$. 
In the counterexample, both fields are the two-element field.)
The finite length property was studied in~\cite{BFKM24}, where it was shown that the equality atoms (Example~\ref{ex:equality-atoms}) and the ordered atoms (Example~\ref{ex:order-atoms}) have this property over any field. 

The main contribution of this paper is to establish the finite length property for more structures. 
We will use two different techniques for this purpose.
